% \documentclass[letterpaper,10pt]{article}

% \usepackage{latexsym}
% \usepackage[empty]{fullpage}
% \usepackage{titlesec}
% \usepackage{marvosym}
% \usepackage[usenames,dvipsnames]{color}
% \usepackage{verbatim}
% \usepackage{enumitem}
% \usepackage[hidelinks]{hyperref}
% \usepackage{fancyhdr}
% \usepackage[english]{babel}
% \usepackage{tabularx}
% \input{glyphtounicode}

% % Font options
% \usepackage{newtxtext,newtxmath}  % Use Times New Roman font

% \pagestyle{fancy}
% \fancyhf{}
% \fancyfoot{}
% \renewcommand{\headrulewidth}{0pt}
% \renewcommand{\footrulewidth}{0pt}

% \addtolength{\oddsidemargin}{-0.5in}
% \addtolength{\evensidemargin}{-0.5in}
% \addtolength{\textwidth}{1in}
% \addtolength{\topmargin}{-0.25in}
% \addtolength{\textheight}{0.5in}

% \urlstyle{same}
% \raggedbottom
% \raggedright
% \setlength{\tabcolsep}{0in}

% % Section formatting
% \titleformat{\section}{\Large\bfseries\scshape\raggedright}{}{0em}{}[\titlerule]

% % Ensure PDF is machine readable
% \pdfgentounicode=1

% % Custom commands
% \newcommand{\resumeItem}[1]{\item\small{#1} \vspace{-3pt}}
% \newcommand{\resumeSubheading}[4]{
% \vspace{-1.0pt}\item
%   \begin{tabular*}{0.97\textwidth}[t]{l@{\extracolsep{\fill}}r}
%     \textbf{#1} & #2 \
%     #3 & #4 \
%   \end{tabular*}\vspace{-7pt}
% }

% \newcommand{\resumeDegreeSubheading}[6]{
% \vspace{-1.0pt}\item
%   \begin{tabular*}{0.97\textwidth}[t]{l@{\extracolsep{\fill}}r}
%     \textbf{#1} & #2 \
%     #3 & #4 \
%     #5 & #6 \
%   \end{tabular*}\vspace{-7pt}
% }


% \renewcommand\labelitemii{$\vcenter{\hbox{\tiny$\bullet$}}$}
% \newcommand{\resumeSubHeadingList}{\begin{itemize}[leftmargin=0.15in, label={}]}
% \newcommand{\resumeSubHeadingListEnd}{\vspace{-6.4pt}\end{itemize}}

% \begin{document}

% \begin{center}
%   \textbf{\Huge Saipranav Venkatakrishnan} \
%   \vspace{0.2em}
%   \large 847-348-0250 $|$ \href{mailto:saipranavvk@gmail.com}{sv34@illinois.edu} $|$ 
%   \href{https://www.linkedin.com/in/saipranav-venkatakrishnan}{linkedin.com/in/saipranav-venkatakrishnan} $|$
%   \href{https://github.com/Sai-Pra}{https://github.com/Sai-Pra}
% \end{center}

% % \section{Summary}
% % Passionate Computer Engineer with a strong background in Low Latency Systems, Operating Systems Development, and GPU Programming. Skilled in C, C++, CUDA C++, Python, and a PyTorch contributor. Passions include 3D Modeling/Design and playing Quadball!
% % Passionate AI/ML engineer with a strong background in deep learning, computer vision, and natural language processing. Skilled in Python, TensorFlow, PyTorch, and various ML libraries. Excellent problem-solving, research, and collaboration abilities. Seeking a challenging role to develop cutting-edge AI solutions.

% % \section{Summary}
% % \vspace{-0.5em}
% % \resumeSubHeadingList
% %   \resumeItem{Senior in Computer Engineering specializing in CPU architecture, RTL design, verification, and operating systems development. Looking for challenging roles in processor and hardware design where I can apply and expand my expertise in high-performance processor systems and system-level software interactions.}
% % \resumeSubHeadingListEnd
% \vspace{-1.5em}
% \section{Education}
% \vspace{-0.5em}
% \resumeSubHeadingList
%   \resumeDegreeSubheading
%       {University of Illinois Urbana-Champaign}{GPA: 3.85}
%       {B.S. in Computer Engineering}{August 2022--May 2026}
%       {John P. Desmond Electrical and Computer Engineering Scholarship Award Recipient}{}
%       % {Expected: M.S.}{August 2026--May 2028}
%       % {}{}
%   \vspace{0.5em}
%   \resumeItem{\textbf{Related Coursework}: Computer Architecture, VLSI Design, Computer Organization and Design, Logic Synthesis, Data Structures and Algorithms, Operating Systems Development, Digital Systems Lab, Applied Parallel Programming, Introduction to Algorithms and Models of Computations}

% \resumeSubHeadingListEnd

% \section{Technical Skills}
% \vspace{-0.5em}
% \resumeSubHeadingList
%   \resumeItem{\textbf{Programming Languages and Skills}: SystemVerilog, C, C++, CUDA C++, Python, Lua, ASIC Design, Computer Vision, Object-Oriented Programming, Parallel and Multithreaded Programming, Networking, SIS, ABC}
%   % \resumeItem{\textbf{Library Contributions}: PyTorch, CME Internal Compute Libraries, }
%   \resumeItem{\textbf{Libraries \& Tools}: Synopsys Design Compiler, Synopsys VCS, Verdi, Verilator, Vivado, Vitis, FPGA, Nvidia Nsight Compute, Nvidia Nsight Systems, Oscilloscope, Git, Arduino, Google Cloud, Docker, Kubernetes, Pytorch, Vulkan, Qiskit}
% \resumeSubHeadingListEnd

% \section{Experience}
% \vspace{-0.5em}
% \resumeSubHeadingList
%     \resumeSubheading
%       {Incoming ASIC Design Intern}{May 2026 -- August 2026}
%       {NVIDIA}{Santa Clara, CA}
%     \resumeSubheading
%       {Undergraduate Hardware Research Assistant}{December 2024 -- Present}
%       {PlatformX}{Champaign, IL}
%       \resumeSubHeadingList
%         % \resumeItem{\textbullet\ Improved \textbf{bandwidth utilization} and \textbf{compute efficiency} of AI Wafer-scale chips by developing architectural simulator of new heterogeneous core designs, achieving up to \textbf{15\% performance gains}}
%         \resumeItem{\textbullet\ Designed an \textbf{architectural simulator} to profile AI Wafer-scale chip designs, helping find improvements to \textbf{bandwidth utilization} and \textbf{compute efficiency} through \textbf{heterogeneous core }designs}
%         \resumeItem{\textbullet\ Extended \textbf{GPGPU-Sim} (\textbf{cycle-accurate simulator}) with \textbf{C++} and automation scripts to enable detailed profiling of \textbf{memory access patterns} and \textbf{interconnect bandwidth}}
%         \resumeItem{\textbullet\ Applied simulator extensions to investigate architectural tradeoffs across various \textbf{high-speed interconnect} designs within \textbf{GPU Processing Clusters (GPCs)}, focusing on interactions between \textbf{L1 cache slices} and \textbf{streaming multiprocessors (SMs)} optimized for \textbf{AI workloads}}
%       \resumeSubHeadingListEnd
%     \vspace{0.2em}
%     \resumeSubheading
%       {Computer Architecture (ECE 411) Course Staff}{August 2025 -- Present}
%       {UIUC Computer Organization and Design}{Champaign, IL}
%       \resumeSubHeadingList
%         % Guided -> Trained (Implies knowledge)
%         \resumeItem{\textbullet\ Trained students in designing and implementing \textbf{out-of-order RISC-V processors} with advanced features such as \textbf{non-blocking pipeline caching}, \textbf{branch prediction} and \textbf{recovery}, out-of-order \textbf{load/store units}, and \textbf{superscalar} execution}
%         \resumeItem{\textbullet\ Mentored students on \textbf{RTL debugging} across \textbf{pipeline} and \textbf{cache} projects, providing targeted design and verification support}
%         \resumeItem{\textbullet\ Orchestrated competition with fellow course staff to profile and test \textbf{30+ RISC-V processors}, ranking them on metrics related to \textbf{Power, Performance, and Area (PPA)}}
%       \resumeSubHeadingListEnd
%     \vspace{0.2em}
%     \resumeSubheading
%       {Graphics Engineering Intern}{May 2025 -- August 2025}
%       {Aechelon Technology}{Roeland Park, KS}
%         \resumeSubHeadingList
%           \resumeItem{\textbullet\ Integrating SMPTE ST 2110-compliant \textbf{UDP transmission} into the pC-Nova platform via NVIDIA Rivermax SDK, utilizing \textbf{GPUDirect and RDMA} with full kernel bypass to eliminate \textbf{CPU bottlenecks}}
%           \resumeItem{\textbullet\ Implemented 5 \textbf{GPU kernels} and custom \textbf{fragment shaders} to transform output media into various color (e.g., YUV, sRGB) and vector spaces for compatibility with heterogeneous display systems}
%           \resumeItem{\textbullet\ Optimizing real-time media streaming by collaborating with senior engineers to leverage hardware-accelerated networking and GPU-based processing}
%     \resumeSubHeadingListEnd
    
%   % \resumeSubheading
%   %     {Software Development Intern}{May 2024--November 2024}
%   %     {CME Group}{Chicago, IL}
%   %     \resumeSubHeadingList
%   %         \resumeItem{\textbullet\ Lead a team of 4 interns in developing an application to \textbf{monitor project deployments} to on-premises and cloud server for any \textbf{abnormal resource allocations} and stalls}
%   %         \resumeItem{\textbullet\ Won \textbf{first place} in an internal \textbf{CME Hackathon} to develop a \textbf{low-latency trading algorithm}, earning $1000 as the prize}
%   %     \resumeSubHeadingListEnd
%   % \resumeSubheading
%   %     {Software Development Intern}{March 2023--August 2023}
%   %     {CME Group}{Champaign, IL}
%   %     \resumeSubHeadingList
%   %         \resumeItem{\textbullet\ Supervised a team of 5 intern through \textbf{creating} an internal \textbf{C++ library} to validate research papers about methods to foretell stock prices through corresponding future market data}
%   %         \resumeItem{\textbullet\ Configured a CME security compliant remote development environment and integrated over \textbf{3000+ lines} into the CME codebase}
%   %     \resumeSubHeadingListEnd
% \resumeSubHeadingListEnd

% \section{Projects}
% \vspace{-0.5em}
% \resumeSubHeadingList
%     \resumeSubheading
%       {RV-32IM RISC-V Processor}{January 2025 -- May 2025} 
%       {Computer Architecture Project}{SystemVerilog, C++}
%       \resumeSubHeadingList
%         \resumeItem{\textbullet\ Designed and verified a pipelined, out-of-order, \textbf{N-way superscalar processor} based on the \textbf{RV32IM ISA}, with explicit register renaming}
%         \resumeItem{\textbullet\ Implemented a split Load/Store Unit for \textbf{speculative memory operations}, branch predictors, and \textbf{non-blocking pipelined cache}, including \textbf{functional coverage} metrics to assess verification completeness}
%         \resumeItem{\textbullet\ Debugged RTL with waveform tracing using \textbf{Verilator} and \textbf{Verdi}; verification involved writing unit-level \textbf{SystemVerilog testbenches} and assertions to validate pipeline stages, and simulations were performed with industry tools (VCS-compatible)}
%         \resumeItem{\textbullet\ Performed synthesis and power-performance-timing (\textbf{PPA}) analysis using \textbf{Synopsys Design Compiler}; evaluated across multiple testcases for critical path delays, area constraints, and \textbf{dynamic power} efficiency}
%       \resumeSubHeadingListEnd
%       \vspace{0.2em}
    
%   % \resumeSubheading
%   %     {Graphics Engine}{May 2024 -- August 2024} 
%   %     {Graphics Project}{C++, Vulkan}
%   %   \resumeSubHeadingList
%   %     \resumeItem{\textbullet\ Building a \textbf{graphics pipeline} from scratch using the \textbf{Vulkan SDK}, with over \textbf{6K lines of C++ code} written across the entire codebase for low-latency rendering}
%   %     \resumeItem{\textbullet\ Implemented 5+ core \textbf{graphics pipeline stages}, including input assembly, rasterization, fragment shading, and framebuffer presentation}
%   %     \resumeItem{\textbullet\ Developed a system to automate Vulkan pipeline creation and manage 5+ \textbf{shader modules} across multiple \textbf{render passes}}
%   %   \resumeSubHeadingListEnd

%     % \resumeSubheading
%     %   {Spiking Neural Network}{October 2024 -- January 2025} 
%     %   {Hardware Synthesis and Neuromorphic Computing}{SystemVerilog, FPGA}
%     % \resumeSubHeadingList
%     %   \resumeItem{\textbullet\ Researched \textbf{neuromorphic architectures} and deployed a \textbf{spiking neural network} on an FPGA trained on the \textbf{Berkeley DROID Dataset} for Robotic Manipulation}
%     %   \resumeItem{\textbullet\ Trained Leaky-Integrate-and-Fire (LIF) neuron weights in \textbf{PyTorch}, performed \textbf{quantization} for hardware compatibility, and mapped parameters to fixed-point representations}
%     %   \resumeItem{\textbullet\ Designed LIF neuron modules using shift-register counters to emulate membrane potential dynamics and developed inter-layer \textbf{FIFO queues} to stage \textbf{asynchronous spike traffic}}
%     %   \resumeItem{\textbullet\ Streamed multi-gigabyte testing data through the \textbf{UART protocol} to the \textbf{FPGA} and back, testing inference accuracy against PyTorch model}
%     % \resumeSubHeadingListEnd
    

%     % \resumeSubheading
%     %   {Audio Visualizer Theremin}{October 2024 -- December 2024} 
%     %   {Hardware Synthesis and Verification}{SystemVerilog, FPGA}
%     %   \resumeSubHeadingList
%     %     \resumeItem{\textbullet\ Designed and implemented an FPGA-based SoC on Xilinx for real-time audio synthesis and VGA visualization, controlled via UART input; integrated a MicroBlaze soft processor with custom hardware modules (audio generator, VGA controller) for real-time processing}
%     %     \resumeItem{\textbullet\ Developed a modular SystemVerilog testbench framework for both component and SoC-level validation, featuring automated stimulus generation, self-checking monitors, and coverage tracking}
%     %     \resumeItem{\textbullet\ Validated and debugged designs in Vivado using assertions and waveform analysis to ensure correctness across audio processing states}
        
%     %   \resumeSubHeadingListEnd
    
%   % \resumeSubheading
%   %     {32-bit Operating System}{January 2024 -- May 2024}
%   %     {Systems Project}{C, x86, QEMU}
%   %     \resumeSubHeadingList
%   %         \resumeItem{\textbullet\ Programmed a 32-bit Operating System with a \textbf{3-terminal interface} for \textbf{IA-32} architecture using \textbf{x86 assembly and C}}
%   %         \resumeItem{\textbullet\ Developed a 2-Radix \textbf{Paging System} and \textbf{Allocator} for both 4KB and 4MB page}
%   %         \resumeItem{\textbullet\ Constructed a simple \textbf{virtual file-system} and over 10 \textbf{system call handlers} to interface with the kernel and the physical file-system}
%   %         \resumeItem{\textbullet\ Designed a \textbf{Round Robin Scheduling algorithm} for switching between various userspace programs}
%   %     \resumeSubHeadingListEnd
%   % \resumeSubheading
%   %     {Convolutional Neural Network Optimization}{January 2024 -- May 2024}
%   %     {Parallel Algorithm and Profiling Project}{CUDA, C++}
%   %     \resumeSubHeadingList
%   %         \resumeItem{\textbullet\ Researched and implemented a scalable \textbf{CNN forward pass layer} using \textbf{CUDA}, and verified optimization efficiency through the use of Nvidia Nsight Systems and Compute}
%   %         \resumeItem{\textbullet\ Optimized convolutions through the use of streams, \textbf{shared memory}, constant memory, matrix unrolling, \textbf{tensor cores}, and using custom \textbf{half2} data types}
%   %     \resumeSubHeadingListEnd
%   % \resumeSubheading
%   %     {Pytorch Contributions}{}
%   %     {AI/ML Contributions}{}
%   %     \resumeSubHeadingList
%   %         \resumeItem{\textbullet\ Contributed to PyTorch by writing documentation during PyTorch annual Docathon}
%   %         \resumeItem{\textbullet\ Connected with PyTorch developers to understand code and write correct and detailed comments about their purpose}
%   %     \resumeSubHeadingListEnd
% \resumeSubHeadingListEnd


% % \section{Student Organizations}
% % \vspace{-0.5em}
% % \resumeSubHeadingList
% %     \resumeSubheading
% %       {Engineering Tutor}{January 2024--Present}
% %       {Eta Kappa Nu Honor Society}{Champaign, IL}
% %       \resumeSubHeadingList
% %       \resumeItem{\textbullet\ Inducted as a \textbf{top 20\% student} in the department; tutor peers in topics ranging from LC-3 \textbf{assembly} and \textbf{C} programming to analog \textbf{signal processing}, \textbf{multivariable calculus}, and \textbf{architecture design}}
% %       \resumeItem{\textbullet\ Create custom study guides and \textbf{host review sessions} across 5+ engineering courses, supporting \textbf{100+ students} in exam preparation}
% %       \resumeSubHeadingListEnd
% %       \vspace{0.2em}

% %     \resumeSubheading
% %         {SIGArch}{Champaign, IL}
% %         {Member}{January 2023 -- Present}
% %         \resumeSubHeadingList
% %           \resumeItem{\textbullet\ Facilitated \textbf{technical discussions and reading groups} on computer architecture topics, including modern \textbf{CPUs}, \textbf{GPUs}, and \textbf{emerging accelerator designs}}
% %           \resumeItem{\textbullet\ Helped coordinate architecture-focused talks and workshops featuring graduate students and faculty working on \textbf{microarchitecture, memory systems, and parallel computing}}
% %           \resumeItem{\textbullet\ Collaborated with other \textbf{computing organizations} to host events bridging architecture with \textbf{systems, compilers, and hardware-software co-design}}
% %           \resumeItem{\textbullet\ Supported \textbf{professional and academic networking events} connecting undergraduate students with graduate researchers and industry professionals in computer architecture}
% %         \resumeSubHeadingListEnd
% %           \vspace{0.2em}

% %     \resumeSubheading
% %       {Social Chair}{May 2023 -- May 2024}
% %       {Association for Computing Machinery}{Champaign, IL}
% %         \resumeSubHeadingList
% %           \resumeItem{\textbullet\ Organized and hosted \textbf{50+ social events}, including weekly social hours that connected engineering students from freshmen to graduate levels}
% %           \resumeItem{\textbullet\ \textbf{Expanded organization membership} by over 300 students through active community engagement and \textbf{cross-disciplinary outreach}}
% %           \resumeItem{\textbullet\ Collaborated with \textbf{4+ campus organizations} to \textbf{plan minority-inclusive events} and co-host a large-scale joint school dance}
% %           \resumeItem{\textbullet\ Assisted in coordinating \textbf{12+ professional networking events}, enabling students to engage with \textbf{industry professionals} across diverse computing disciplines}
% %     \resumeSubHeadingListEnd
    


% % \resumeSubHeadingListEnd
% \end{document}



\documentclass[letterpaper,10pt]{article}

\usepackage{latexsym}
\usepackage[empty]{fullpage}
\usepackage{titlesec}
\usepackage{marvosym}
\usepackage[usenames,dvipsnames]{color}
\usepackage{verbatim}
\usepackage{enumitem}
\usepackage[hidelinks]{hyperref}
\usepackage{fancyhdr}
\usepackage[english]{babel}
\usepackage{tabularx}
\input{glyphtounicode}

\usepackage{newtxtext,newtxmath}

\pagestyle{fancy}
\fancyhf{}
\fancyfoot{}
\renewcommand{\headrulewidth}{0pt}
\renewcommand{\footrulewidth}{0pt}

\addtolength{\oddsidemargin}{-0.5in}
\addtolength{\evensidemargin}{-0.5in}
\addtolength{\textwidth}{1in}
\addtolength{\topmargin}{-0.25in}
\addtolength{\textheight}{0.5in}

\urlstyle{same}
\raggedbottom
\raggedright
\setlength{\tabcolsep}{0in}

\titleformat{\section}{\Large\bfseries\scshape\raggedright}{}{0em}{}[\titlerule]

\pdfgentounicode=1

\newcommand{\resumeItem}[1]{\item\small{#1} \vspace{-3pt}}
\newcommand{\resumeSubheading}[4]{
\vspace{-1.0pt}\item
  \begin{tabular*}{0.97\textwidth}[t]{l@{\extracolsep{\fill}}r}
    \textbf{#1} & #2 \
    #3 & #4 \
  \end{tabular*}\vspace{-7pt}
}
\newcommand{\resumeDegreeSubheading}[6]{
\vspace{-1.0pt}\item
  \begin{tabular*}{0.97\textwidth}[t]{l@{\extracolsep{\fill}}r}
    \textbf{#1} & #2 \
    #3 & #4 \
    #5 & #6 \
  \end{tabular*}\vspace{-7pt}
}

\renewcommand\labelitemii{$\vcenter{\hbox{\tiny$\bullet$}}$}
\newcommand{\resumeSubHeadingList}{\begin{itemize}[leftmargin=0.15in, label={}]}
\newcommand{\resumeSubHeadingListEnd}{\vspace{-6.4pt}\end{itemize}}

\begin{document}

\begin{center}
  \textbf{\Huge Saipranav Venkatakrishnan} \
  \vspace{0.2em}
  \large 847-348-0250 $|$ \href{mailto:saipranavvk@gmail.com}{saipranavvk@gmail.com} $|$
  \href{https://www.linkedin.com/in/saipranav-venkatakrishnan}{linkedin.com/in/saipranav-venkatakrishnan} $|$
  \href{https://github.com/Sai-Pra}{github.com/Sai-Pra}
\end{center}

\vspace{-1.5em}
\section{Education}
\vspace{-0.5em}
\resumeSubHeadingList
  \resumeDegreeSubheading
      {University of Illinois Urbana-Champaign}{}
      {B.S. in Computer Engineering}{August 2022 -- May 2026 (GPA: 3.85)}
      {M.S./PhD in Computer Engineering}{August 2026 -- May 2032}
  \vspace{0.5em}
  \resumeItem{\textbf{Skills \& Tools}: SystemVerilog, Verilog, C/C++, CUDA C++, Python; Synopsys Design Compiler, VCS, Verdi, Verilator, Yosys, Spyglass Linter, CDC/RDC Analysis, Nvidia Nsight Compute/Systems, Git, PyTorch; ASIC \& RTL Design, Hardware--Software Co-design, Computer Architecture, VLSI Design, Logic Synthesis, Parallel Programming}
\resumeSubHeadingListEnd

% \section{Technical Skills}
% \vspace{-0.5em}
% \resumeSubHeadingList
%   \resumeItem{\textbf{Languages \& Design}: SystemVerilog, Verilog, C, C++, CUDA C++, Python, Lua; ASIC Design, RTL Design, Hardware--Software Co-design, Parallel \& Multithreaded Programming}
%   \resumeItem{\textbf{Tools \& Flows}: Synopsys Design Compiler, Synopsys VCS, Verdi, Verilator, Yosys, lint (Spyglass), CDC/RDC Analysis, Vivado, Vitis, FPGA, Nvidia Nsight Compute, Nvidia Nsight Systems, Git, Docker, Kubernetes, PyTorch, Vulkan, Qiskit, SIS, ABC}
% \resumeSubHeadingListEnd

\section{Experience}
\vspace{-0.5em}
\resumeSubHeadingList

    % -----------------------------------------------------------------------
    % NVIDIA -- now with real bullets from NVLink Chip2Chip work
    % CDC is explicitly called out in the OpenAI JD tools/methodology list;
    % packetization + fanout reduction + async/sync transactions are all
    % directly relevant to compute, memory, and interconnect components
    % -----------------------------------------------------------------------
    \resumeSubheading
      {ASIC Design Intern --- NVLink}{May 2026 -- August 2026}
      {NVIDIA}{Santa Clara, CA}
      \resumeSubHeadingList
        \resumeItem{\textbullet\ Designed \textbf{20+ production Verilog modules} for the \textbf{NVLink Chip-to-Chip interconnect}, implementing \textbf{packet framing/deframing}, \textbf{interconnect fanout reduction}, and \textbf{asynchronous/synchronous transaction} handling}
        \resumeItem{\textbullet\ Implemented \textbf{CDC (Clock Domain Crossing)} logic for high-speed chip-to-chip links, applying synchronizer chains and handshake protocols to ensure \textbf{functional correctness} and \textbf{timing closure} across async clock boundaries}
        \resumeItem{\textbullet\ Collaborated with \textbf{DV and Physical Design} teams to verify module integration, validate \textbf{CDC/RDC} constraints, and meet area and timing targets within the full NVLink SoC design flow}
      \resumeSubHeadingListEnd
    \vspace{0.2em}

    % -----------------------------------------------------------------------
    % PlatformX -- re-led with co-design and AI workload language from JD
    % -----------------------------------------------------------------------
    \resumeSubheading
      {Undergraduate AI Hardware Research Assistant}{December 2024 -- Present}
      {PlatformX}{Champaign, IL}
      \resumeSubHeadingList
        \resumeItem{\textbullet\ Helped develop an \textbf{architectural simulator} for \textbf{AI wafer-scale accelerators}, enabling \textbf{hardware--software co-design} studies that identify improvements in \textbf{bandwidth utilization} and \textbf{compute efficiency} through reconfigurable designs}
        \resumeItem{\textbullet\ Modeled \textbf{AI/ML workload dataflow} across accelerator subsystems to evaluate \textbf{performance--area--power tradeoffs}, providing architectural feedback to guide \textbf{compute and memory subsystem} design decisions}
        \resumeItem{\textbullet\ Extended \textbf{GPGPU-Sim} (cycle-accurate simulator) with \textbf{C++} to enable detailed profiling of \textbf{memory access patterns} and \textbf{interconnect bandwidth}; applied across \textbf{GPU Processing Cluster (GPC)} designs studying \textbf{L1 cache slice} and \textbf{SM} interactions optimized for \textbf{AI workloads}}
      \resumeSubHeadingListEnd
    \vspace{0.2em}

    % -----------------------------------------------------------------------
    % ECE 411 Course Staff
    % -----------------------------------------------------------------------
    \resumeSubheading
      {Computer Architecture (ECE 411) Course Staff}{August 2025 -- Present}
      {UIUC Computer Organization and Design}{Champaign, IL}
      \resumeSubHeadingList
        \resumeItem{\textbullet\ Trained students in designing and implementing \textbf{out-of-order RISC-V processors} with \textbf{non-blocking pipeline caching}, \textbf{branch prediction \& recovery}, out-of-order \textbf{load/store units}, and \textbf{superscalar} execution}
        \resumeItem{\textbullet\ Mentored students on \textbf{RTL debugging} across pipeline and cache projects using \textbf{Verdi} and \textbf{Verilator}; orchestrated competition profiling \textbf{30+ RISC-V processors} on \textbf{Power, Performance, and Area (PPA)} metrics}
      \resumeSubHeadingListEnd
    \vspace{0.2em}

    % -----------------------------------------------------------------------
    % Aechelon
    % -----------------------------------------------------------------------
    % \resumeSubheading
    %   {Graphics Engineering Intern}{May 2025 -- August 2025}
    %   {Aechelon Technology}{Roeland Park, KS}
    %   \resumeSubHeadingList
    %     \resumeItem{\textbullet\ Integrated SMPTE ST 2110-compliant \textbf{UDP transmission} into the pC-Nova platform via NVIDIA Rivermax SDK, utilizing \textbf{GPUDirect and RDMA} with full kernel bypass to eliminate \textbf{CPU bottlenecks}}
    %     \resumeItem{\textbullet\ Implemented 5 \textbf{GPU kernels} and custom \textbf{fragment shaders} for color space conversion (YUV, sRGB) across heterogeneous display systems; optimized \textbf{hardware-accelerated} real-time media streaming with senior engineers}
    %   \resumeSubHeadingListEnd

    \resumeSubheading
      {Graphics Engineering Intern}{May 2025 -- August 2025}
      {Aechelon Technology}{Roeland Park, KS}
        \resumeSubHeadingList
          \resumeItem{\textbullet\ Integrating SMPTE ST 2110-compliant \textbf{UDP transmission} into the pC-Nova platform via NVIDIA Rivermax SDK, utilizing \textbf{GPUDirect and RDMA} with full kernel bypass to eliminate \textbf{CPU bottlenecks}}
          \resumeItem{\textbullet\ Implemented 5 \textbf{GPU kernels} and custom \textbf{fragment shaders} to transform output media into various color (e.g., YUV, sRGB) and vector spaces for compatibility with heterogeneous display systems}
          \resumeItem{\textbullet\ Optimizing real-time media streaming by collaborating with senior engineers to leverage hardware-accelerated networking and GPU-based processing}
    \resumeSubHeadingListEnd
    
  \resumeSubheading
      {Software Development Intern}{May 2024--November 2024}
      {CME Group}{Chicago, IL}
      \resumeSubHeadingList
          \resumeItem{\textbullet\ Lead a team of 4 interns in developing an application to \textbf{monitor project deployments} to on-premises and cloud server for any \textbf{abnormal resource allocations} and stalls}
          \resumeItem{\textbullet\ Won \textbf{first place} in an internal \textbf{CME Hackathon} to develop a \textbf{low-latency trading algorithm}, earning $1000 as the prize}
      \resumeSubHeadingListEnd
  \resumeSubheading
      {Software Development Intern}{March 2023--August 2023}
      {CME Group}{Champaign, IL}
      \resumeSubHeadingList
          \resumeItem{\textbullet\ Supervised a team of 5 intern through \textbf{creating} an internal \textbf{C++ library} to validate research papers about methods to foretell stock prices through corresponding future market data}
          \resumeItem{\textbullet\ Configured a CME security compliant remote development environment and integrated over \textbf{3000+ lines} into the CME codebase}
      \resumeSubHeadingListEnd
\resumeSubHeadingListEnd


\section{Projects}
\vspace{-0.5em}
\resumeSubHeadingList

    % -----------------------------------------------------------------------
    % FORTE -- now with rich detail from the paper
    % Tape-out track record is the #1 thing the JD asks for.
    % Lead with the scope, then ECC/scrubbing, then RoT/Keystone co-design,
    % then synthesis/area numbers (concrete = credible)
    % -----------------------------------------------------------------------
    \resumeSubheading
      {FORTE: Fault-Tolerant RISC-V Processor Tape-Out}{April 2025 -- Present}
      {Silicon Research Project (65nm Tape-Out)}{SystemVerilog, VCS Design Compiler, Verdi, Spyglass, Yosys}
      \resumeSubHeadingList
        \resumeItem{\textbullet\ Designed and implemented \textbf{fault-tolerance RTL} on a \textbf{CORE-V Wally} (RV64IMA, Sv39 MMU, PMP, M/S-mode) baseline targeting \textbf{65nm tape-out}; features include \textbf{SECDED ECC} across register file, L1 I/D-cache, PTW/TLB, a \textbf{cacheline and CSR scrubbing FSM}, and a \textbf{shadow pipeline} for redundant execution-unit fault detection}
        \resumeItem{\textbullet\ Implemented \textbf{clock integrity monitor} (reference-oscillator comparison, $\pm\% glitch threshold, pipeline flush/reset response) and \textbf{micro-TMR hardening} for security-critical CSRs (\texttt{mstatus}, \texttt{mepc}, \texttt{mtvec}, PMP regs) to defend against physical fault-injection attacks}
        \resumeItem{\textbullet\ Built complete \textbf{hardware Root of Trust} infrastructure: boot ROM with \textbf{SHA-256/RSA-2048} image authentication, \textbf{TRNG} via prime-length ring oscillators (Zkr), and full Keystone \textbf{TEE co-design} (16 PMP CSRs, architectural state scrubbing on enclave entry/exit, debug port lockdown)}
        \resumeItem{\textbullet\ Conducted area feasibility study using \textbf{VCS Design Compiler} using TSMC's 65nm PDK; designed \textbf{AXI-Lite chip-to-board interface} ($\sim0 pins) for FPGA-based bringup}
      \resumeSubHeadingListEnd
      \vspace{0.2em}

    % -----------------------------------------------------------------------
    % RV-32IM OOO Processor
    % -----------------------------------------------------------------------
    \resumeSubheading
      {RV-32IM Out-of-Order RISC-V Processor}{January 2025 -- May 2025}
      {Computer Architecture Project}{SystemVerilog, C++}
      \resumeSubHeadingList
        \resumeItem{\textbullet\ Designed and verified a pipelined, out-of-order, \textbf{N-way superscalar processor} based on the \textbf{RV32IM ISA}, with explicit register renaming}
        \resumeItem{\textbullet\ Implemented a split Load/Store Unit for \textbf{speculative memory operations}, branch predictors, and \textbf{non-blocking pipelined cache}, including \textbf{functional coverage} metrics to assess verification completeness}
        \resumeItem{\textbullet\ Debugged RTL with waveform tracing using \textbf{Verilator} and \textbf{Verdi}; verification involved writing unit-level \textbf{SystemVerilog testbenches} and assertions to validate pipeline stages, and simulations were performed with industry tools (VCS-compatible)}
        \resumeItem{\textbullet\ Performed synthesis and power-performance-timing (\textbf{PPA}) analysis using \textbf{Synopsys Design Compiler}; evaluated across multiple testcases for critical path delays, area constraints, and \textbf{dynamic power} efficiency}
      \resumeSubHeadingListEnd
      \vspace{0.2em}

  \resumeSubheading
      {Graphics Engine}{May 2024 -- August 2024} 
      {Graphics Project}{C++, Vulkan}
    \resumeSubHeadingList
      \resumeItem{\textbullet\ Building a \textbf{graphics pipeline} from scratch using the \textbf{Vulkan SDK}, with over \textbf{6K lines of C++ code} written across the entire codebase for low-latency rendering}
      \resumeItem{\textbullet\ Implemented 5+ core \textbf{graphics pipeline stages}, including input assembly, rasterization, fragment shading, and framebuffer presentation}
      \resumeItem{\textbullet\ Developed a system to automate Vulkan pipeline creation and manage 5+ \textbf{shader modules} across multiple \textbf{render passes}}
    \resumeSubHeadingListEnd

    \resumeSubheading
      {Spiking Neural Network}{October 2024 -- January 2025} 
      {Hardware Synthesis and Neuromorphic Computing}{SystemVerilog, FPGA}
    \resumeSubHeadingList
      \resumeItem{\textbullet\ Researched \textbf{neuromorphic architectures} and deployed a \textbf{spiking neural network} on an FPGA trained on the \textbf{Berkeley DROID Dataset} for Robotic Manipulation}
      \resumeItem{\textbullet\ Trained Leaky-Integrate-and-Fire (LIF) neuron weights in \textbf{PyTorch}, performed \textbf{quantization} for hardware compatibility, and mapped parameters to fixed-point representations}
      \resumeItem{\textbullet\ Designed LIF neuron modules using shift-register counters to emulate membrane potential dynamics and developed inter-layer \textbf{FIFO queues} to stage \textbf{asynchronous spike traffic}}
      \resumeItem{\textbullet\ Streamed multi-gigabyte testing data through the \textbf{UART protocol} to the \textbf{FPGA} and back, testing inference accuracy against PyTorch model}
    \resumeSubHeadingListEnd
    

    \resumeSubheading
      {Audio Visualizer Theremin}{October 2024 -- December 2024} 
      {Hardware Synthesis and Verification}{SystemVerilog, FPGA}
      \resumeSubHeadingList
        \resumeItem{\textbullet\ Designed and implemented an FPGA-based SoC on Xilinx for real-time audio synthesis and VGA visualization, controlled via UART input; integrated a MicroBlaze soft processor with custom hardware modules (audio generator, VGA controller) for real-time processing}
        \resumeItem{\textbullet\ Developed a modular SystemVerilog testbench framework for both component and SoC-level validation, featuring automated stimulus generation, self-checking monitors, and coverage tracking}
        \resumeItem{\textbullet\ Validated and debugged designs in Vivado using assertions and waveform analysis to ensure correctness across audio processing states}
        
      \resumeSubHeadingListEnd
    
  \resumeSubheading
      {32-bit Operating System}{January 2024 -- May 2024}
      {Systems Project}{C, x86, QEMU}
      \resumeSubHeadingList
          \resumeItem{\textbullet\ Programmed a 32-bit Operating System with a \textbf{3-terminal interface} for \textbf{IA-32} architecture using \textbf{x86 assembly and C}}
          \resumeItem{\textbullet\ Developed a 2-Radix \textbf{Paging System} and \textbf{Allocator} for both 4KB and 4MB page}
          \resumeItem{\textbullet\ Constructed a simple \textbf{virtual file-system} and over 10 \textbf{system call handlers} to interface with the kernel and the physical file-system}
          \resumeItem{\textbullet\ Designed a \textbf{Round Robin Scheduling algorithm} for switching between various userspace programs}
      \resumeSubHeadingListEnd
  \resumeSubheading
      {Convolutional Neural Network Optimization}{January 2024 -- May 2024}
      {Parallel Algorithm and Profiling Project}{CUDA, C++}
      \resumeSubHeadingList
          \resumeItem{\textbullet\ Researched and implemented a scalable \textbf{CNN forward pass layer} using \textbf{CUDA}, and verified optimization efficiency through the use of Nvidia Nsight Systems and Compute}
          \resumeItem{\textbullet\ Optimized convolutions through the use of streams, \textbf{shared memory}, constant memory, matrix unrolling, \textbf{tensor cores}, and using custom \textbf{half2} data types}
      \resumeSubHeadingListEnd
  \resumeSubheading
      {Pytorch Contributions}{}
      {AI/ML Contributions}{}
      \resumeSubHeadingList
          \resumeItem{\textbullet\ Contributed to PyTorch by writing documentation during PyTorch annual Docathon}
          \resumeItem{\textbullet\ Connected with PyTorch developers to understand code and write correct and detailed comments about their purpose}
      \resumeSubHeadingListEnd
\resumeSubHeadingListEnd


\section{Student Organizations}
\vspace{-0.5em}
\resumeSubHeadingList
    \resumeSubheading
      {Engineering Tutor}{January 2024--Present}
      {Eta Kappa Nu Honor Society}{Champaign, IL}
      \resumeSubHeadingList
      \resumeItem{\textbullet\ Inducted as a \textbf{top 20\% student} in the department; tutor peers in topics ranging from LC-3 \textbf{assembly} and \textbf{C} programming to analog \textbf{signal processing}, \textbf{multivariable calculus}, and \textbf{architecture design}}
      \resumeItem{\textbullet\ Create custom study guides and \textbf{host review sessions} across 5+ engineering courses, supporting \textbf{100+ students} in exam preparation}
      \resumeSubHeadingListEnd
      \vspace{0.2em}

    \resumeSubheading
        {SIGArch}{Champaign, IL}
        {Member}{January 2023 -- Present}
        \resumeSubHeadingList
          \resumeItem{\textbullet\ Facilitated \textbf{technical discussions and reading groups} on computer architecture topics, including modern \textbf{CPUs}, \textbf{GPUs}, and \textbf{emerging accelerator designs}}
          \resumeItem{\textbullet\ Helped coordinate architecture-focused talks and workshops featuring graduate students and faculty working on \textbf{microarchitecture, memory systems, and parallel computing}}
          \resumeItem{\textbullet\ Collaborated with other \textbf{computing organizations} to host events bridging architecture with \textbf{systems, compilers, and hardware-software co-design}}
          \resumeItem{\textbullet\ Supported \textbf{professional and academic networking events} connecting undergraduate students with graduate researchers and industry professionals in computer architecture}
        \resumeSubHeadingListEnd
          \vspace{0.2em}

    \resumeSubheading
      {Social Chair}{May 2023 -- May 2024}
      {Association for Computing Machinery}{Champaign, IL}
        \resumeSubHeadingList
          \resumeItem{\textbullet\ Organized and hosted \textbf{50+ social events}, including weekly social hours that connected engineering students from freshmen to graduate levels}
          \resumeItem{\textbullet\ \textbf{Expanded organization membership} by over 300 students through active community engagement and \textbf{cross-disciplinary outreach}}
          \resumeItem{\textbullet\ Collaborated with \textbf{4+ campus organizations} to \textbf{plan minority-inclusive events} and co-host a large-scale joint school dance}
          \resumeItem{\textbullet\ Assisted in coordinating \textbf{12+ professional networking events}, enabling students to engage with \textbf{industry professionals} across diverse computing disciplines}
    \resumeSubHeadingListEnd
    

\resumeSubHeadingListEnd

\end{document}
